\part{物理问题}

\chapter{多体问题}

\section{质心系}


\section{弹性碰撞}

\section{Virial定理}


\chapter{中心势}
\label{sec:celestial_pot}
运动质点所受的力的作用线始终通过惯性系中一固定点力心，即为中心力。

中心力大小为力心到质点的距离$r$的函数，方向始终沿力心到质点的连线，即
\begin{equation}
    \vec{F} = f(r)\hat{\vec{r}},
\end{equation}
则中心力为保守力，可以用势能来描述$V(r)$，则
\begin{equation}
    \lag = \frac12 m \dot{\vec{r}}^2 - V(r)
\end{equation}


\section{Kepler问题}
\begin{equation}
     V(r)\propto \frac{1}{r} = -\frac{\alpha}{r}
\end{equation}
那么我们可以在极坐标系下讨论该问题，即Lagrangian可以写为
\begin{equation}
    \lag = \frac{1}{2}m(\dot r ^2 + r^2\dot \phi^2) + \frac{\alpha}{r}
\end{equation}

\section{轨道方程与进动}

\section{轨道稳定性}

\section{Lotherford散射}

\chapter{正则微扰论}
\label{sec:perturbation}
\section{含时微扰论}

\section{定态微扰论}

\section{久期方程}
\section{浸渐不变量}


\chapter{转动}
\label{sec:rotation}

\section{三维转动}

\section{转动惯量}

\section{Euler定理}

\section{转动参考系}

\chapter{微振动}
\label{sec:vibration}

\section{一维自由振动}

\section{强迫振动与阻尼振动}



\section{多维简谐振动}

\section{连续介质力学}
\label{sec:continuity field}

\section{应力与应变张量}

\chapter{混沌}
\label{sec:chaos}

\section{三体问题}

\section{参变共振}

\section{周期运动与Poincar\'e映射}

\section{Logistic方程}